本作品基于安路 EG4S20 FPGA (HX4S20C 开发板) 设计并实现了“赛博美学数字艺术多媒体交互画框”，全面完成赛题的全部基础要求与五类拓展功能。系统以一套固化镜像集成 TF 卡 BMP 图片读取、HDMI 720P 音视频输出、手动切图、可配置周期的自动轮播和上电自启动，并将滚动字幕、淡出淡入转场、多尺寸图片等比缩放、参数实时调节与真实音频可视化融入统一的展示界面，使静态素材播放与动态数字艺术在同一设备上协同运行。

作品的核心创新是“四季动态夜城”，将赛博美学中的霓虹色彩、城市层次与天气变化组织为雨夜霓虹、黄昏雷云、薄雾深夜和飘雪清晨四种氛围。画面不依赖预录视频或连续读取动画帧。画面不依赖预录视频，而由FPGA根据像素坐标、季节色板和时间状态逐像素实时生成：雨雪、云雾、灯带和招牌持续运动，街区广告屏随正在播放的音乐跳动。通过板载按键、拨码开关和中文菜单即可切图、选曲、切换季节，实时调整亮度、对比度与音量；字幕随展示内容自动更新。

系统采用多时钟域并行流水架构，4 个 PLL 分别为 TF 访问、SDRAM 存储、画面合成和 HDMI 输出提供时钟。SDRAM 六帧缓存使切图无需重新读卡，完整帧发布与帧边界切换配合淡出淡入转场，保持画面呈现的连续性。内部画布扫描后输出 $1280\times720$、60 Hz 视频，音频经同一 HDMI 接口同步传输。设计在 EG4S20 上使用 84\% 的 LUT 和全部 32Kb RAM 块，各时钟域均无时序违例。

作品无需上位机，素材准备完成后即可离线独立运行。该设计将图片播放、动态场景生成与实时交互整合为可操作的数字艺术画框，适用于数字广告牌、商业展陈终端、应急信息发布、车载辅助显示及教育演示等需要长期展示音画内容的场合。

\noindent\textbf{关键词：} \mbox{720P}；\mbox{HDMI}；\mbox{赛博美学}；\mbox{数字艺术画框}；\mbox{动态场景生成}；\mbox{实时交互}；\mbox{音画同步}；\mbox{图像缩放}；\mbox{音频可视化}
